% 对应题号 M1-07；图与本题条件、推导、结论逐项对应。
% 图契约：印刷宽16cm；最小字号8.4pt；几何/场或曲线在左，计算关系在右。
\documentclass[tikz,border=7pt]{standalone}
\usepackage{fontspec,xeCJK,amsmath,amssymb,graphicx}
\IfFontExistsTF{Times New Roman}{\setmainfont{Times New Roman}}{\setmainfont{TeX Gyre Termes}}
\IfFontExistsTF{Heiti SC}{\setCJKmainfont[AutoFakeBold=1.5]{Heiti SC}}{\setCJKmainfont[AutoFakeBold=1.5]{FandolHei-Regular}}
\usetikzlibrary{arrows.meta,bending,calc,patterns,positioning}
\definecolor{ink}{HTML}{243745}
\definecolor{blue}{HTML}{275DAB}
\definecolor{green}{HTML}{177C63}
\definecolor{red}{HTML}{BD4935}
\definecolor{orange}{HTML}{AA6C16}
\definecolor{purple}{HTML}{715393}
\tikzset{every node/.style={font=\fontsize{9.3}{12.5}\selectfont,text=ink},
 axis/.style={-{Stealth[length=4pt]},line width=.65pt,draw=ink!65},
 vec/.style={-{Stealth[length=5pt]},line width=1pt,draw=blue},
 short/.style={-{Stealth[length=3pt]},line width=.8pt},
 guide/.style={dashed,draw=ink!40,line width=.5pt},
 card/.style={rounded corners=2pt,draw=blue!65,fill=blue!4,align=center,text width=6.6cm,inner xsep=9pt,inner ysep=7pt},
 title/.style={font=\fontsize{10.5}{14}\selectfont\bfseries,text width=15cm,align=center},
 label/.style={font=\fontsize{8.4}{11}\selectfont,align=center}}
\begin{document}
\begin{tikzpicture}[x=1cm,y=1cm]
\path[use as bounding box] (-8,-3.4) rectangle(8,3.4);
\node[title] at (0,3) {M1-07\quad 同一点沿不同坐标方向的变化率};
\begin{scope}[shift={(-4,0)}]

\draw[axis] (-2.3,.65)--(2.3,.65) node[right] {$x$};\draw[axis] (0,.45)--(0,2.15) node[above] {$f$};
\draw[blue,line width=1pt] plot[domain=-1.6:1.6,samples=60] (\x,{.65+.35*\x*\x});
\node[label,blue] at (-1.45,2.05) {$y,z$ 固定};
\draw[axis] (-2.3,-1.65)--(2.3,-1.65) node[right] {$y$};\draw[axis] (0,-2.5)--(0,-.2) node[above] {$f$};
\draw[green,line width=1pt] plot[domain=-1.3:1.9] (\x,{-1.65+.55*\x});
\node[label,green] at (-1.4,-.45) {$x,z$ 固定};
\node[label] at (0,-2.85) {分别示意正参数切片；求导公式对完整场成立};

\end{scope}\node[card,draw=blue!65,fill=blue!4] at (3.95,1.65) {$f=x^2y+z$};
\node[card,draw=green!65,fill=green!4] at (3.95,0) {$\displaystyle\partial_x f=2xy$\\上图读 $x$ 切片的局部斜率。};
\node[card,draw=orange!65,fill=orange!4] at (3.95,-1.65) {$\displaystyle\partial_y f=x^2$\\下图是关于 $y$ 的一次函数。};
\end{tikzpicture}\end{document}
